% =====================================================================
\section{Part III, AI Exercise 2: Predicting Cross-Sectional Deviations}
% =====================================================================

We asked an LLM which of the 49 industries it expected to fall above, below and
near the security market line, and recorded its answer before comparing it with
our results from Question (d).

\paragraph{(a) The LLM's predictions.}
\begin{tight}
\item \textit{Above the line:} Tobacco, Utilities, Beer \& Liquor / Food.
\item \textit{Below the line:} Coal / Mining, Computers / Electronic Equipment
/ Semiconductors, Autos / Steel.
\item \textit{Near the line:} Wholesale / Retail, Transportation, Financials /
Business Services.
\end{tight}

Its reasoning was the low-beta anomaly. Defensive industries have low beta but
still earn competitive returns, so they plot above the line, while cyclical
industries have high beta without the returns to match. It added a sin-stock
premium for tobacco and a glamour effect for technology. It stated that these
were predictions from documented patterns, not calculations from the data.

\paragraph{(b) What the plot shows.} Deviations are measured from the
theoretical SML in Figure~\ref{fig:sml}, as
$\overline{R_i} - \bar{R}_M^{e}\,b_{iM}$.

\begin{table}[H]
\centering
\caption{Largest deviations from the theoretical security market line, balanced
sample, percent per month. A positive value means the industry earned more than
its beta implies.}
\small
\begin{tabular}{lrlrlr}
\toprule
\multicolumn{2}{c}{Above} & \multicolumn{2}{c}{Below}
& \multicolumn{2}{c}{Near} \\
\cmidrule(lr){1-2}\cmidrule(lr){3-4}\cmidrule(lr){5-6}
Smoke & $+0.675$ & Other & $-0.468$ & BldMt & $-0.003$ \\
Guns  & $+0.461$ & RlEst & $-0.429$ & Boxes & $+0.004$ \\
Soda  & $+0.274$ & PerSv & $-0.376$ & Telcm & $+0.006$ \\
\bottomrule
\end{tabular}
\label{tab:dev}
\end{table}

\paragraph{(c) How accurate were the predictions?} The above-line predictions
were accurate. All three industries plot above the line, and Tobacco is the
largest positive deviation in the sample. The near-line predictions were also
reasonable: Wholesale, Transportation and Financials all sit close to the line,
though none is among the three closest.

The below-line predictions were mostly wrong. Coal and Mining plot slightly
above the line rather than below it. Of the technology group, only Software is
clearly below, while Chips and Electronic Equipment are above. Only Steel was
placed correctly, and none of the three largest negative deviations was
anticipated.

This asymmetry is explained by how the model reasoned. It used beta alone,
predicting that low-beta industries plot above the line and high-beta
industries plot below. Because the estimated line is so flat, the first half of
this is almost automatic, since nearly any low-beta industry lands above the
theoretical SML. That is why the above-line predictions worked. The below-line
predictions needed more, because high beta is not sufficient: Chips and
Electronic Equipment have high betas but earned enough to stay above the line.
The industries that actually fall furthest below have betas near the sample
average and unusually low returns, which reflects industry-specific weakness
rather than a risk story the model could infer.

Its caveat about technology was partly right. It predicted technology would
fall below the line but noted that recent data might reverse this. The reversal
is what we observe for Chips and Electronic Equipment, though not for Software.
